% Generated draft: requires mathematical and source review.
\documentclass[11pt]{article}
\usepackage[margin=25mm]{geometry}
\usepackage{fontspec}
\setmainfont{Latin Modern Roman}
\usepackage{amsmath,amssymb,mathtools,unicode-math}
\setmathfont{Latin Modern Math}
\usepackage{enumitem,needspace,xurl,hyperref}
\setlength{\parindent}{0pt}
\setlength{\parskip}{5pt}
\setlength{\emergencystretch}{3em}
\newcommand{\N}{\mathbb{N}}
\newcommand{\Z}{\mathbb{Z}}
\newcommand{\Q}{\mathbb{Q}}
\newcommand{\R}{\mathbb{R}}
\newcommand{\C}{\mathbb{C}}
\newcommand{\F}{\mathbb{F}}
\newcommand{\A}{\mathbb{A}}
\newcommand{\PP}{\mathbb{P}}
\newcommand{\E}{\mathbb{E}}
\newcommand{\eps}{\varepsilon}
\newcommand{\dd}{\,\mathrm{d}}
\begin{document}
\section*{UnsolvedMath ID: 11000062}
\textbf{Source problem code:} catalogue entry
\subsection{Definitions and mathematical statement}
Let \(\operatorname{Mod}_g^{\pm}\) be the extended mapping class group, including orientation-reversing classes, and let \(H_k=\mathcal I_g(k)\) be the source's filtration subgroup for \(k\ge1\). Its exact definition and genus hypotheses are not provided. The summary calls these Johnson-filtration terms, while the closing sentence refers to a lower central series; those filtrations are not silently identified.

The automorphism group \(\operatorname{Aut}(H_k)\) consists of all abstract automorphisms of \(H_k\). Its abstract commensurator \(\operatorname{Comm}(H_k)\) consists of isomorphisms between finite-index subgroups, identified when they agree on a further finite-index subgroup. For a filtration invariant under \(\operatorname{Mod}_g^{\pm}\), conjugation gives natural maps
\[
 \operatorname{Mod}_g^{\pm}\longrightarrow\operatorname{Aut}(H_k)
 \quad\text{and}\quad
 \operatorname{Mod}_g^{\pm}\longrightarrow\operatorname{Comm}(H_k).
\]
Question 5.18 asks for which \(k\ge1\) the source's equalities
\[
 \operatorname{Aut}(\mathcal I_g(k))=\operatorname{Mod}_g^{\pm},
 \qquad
 \operatorname{Comm}(\mathcal I_g(k))=\operatorname{Mod}_g^{\pm}
\]
hold. A usual rigidity interpretation is that the respective natural conjugation maps are isomorphisms; literal set equality is not intended. The source does not explicitly distinguish that interpretation from abstract group isomorphism.

The source says Theorem 5.17 answers the questions for \(k=1,2\), but the theorem's statement and hypotheses are missing, so they are not reconstructed here. Its motivation is recovery of the ambient mapping class group from deeper subgroups. The summary reports known Torelli/filtration automorphism results, mentioning genus at least \(3\), but no all-\(k\) result or complete references. The two requested rigidity properties are retained separately.
\subsection{Short English statement}
Determine which filtration terms have full automorphism group and abstract commensurator equal, in the intended natural sense, to the extended mapping class group. The source cites answers at the first two terms, but the exact filtration and genus hypotheses are missing.
\subsection{Sources}
\begin{itemize}
\item Local UnsolvedMath registry, permanent ID \(11000062\), exact source code \texttt{catalogue entry}, Question 5.18, PDF page 50. No individual proposer or proposed year is supplied. The statement cross-refers to Theorem 5.17 for \(k=1,2\).
\item Farb, edited volume \emph{Problems on Mapping Class Groups and Related Topics}: \url{https://www.math.uchicago.edu/~farb/papers/mcgbook.pdf}. The volume has 329 explicitly typeset numbered items across contributed chapters, repeated numbering distinguished by PDF page.
\item Extraction pdf-font-structure, accessed 2026-07-29. Source evidence NEEDS\_REVIEW presents an open problem and calls for a current primary-literature audit. Supplied OPEN-TRIAGE summary mentions Torelli/Johnson-filtration automorphism work and genus-at-least-3 results, but gives no precise theorem or full citations and says all \(k\) are not settled.
\item Dataset attribution: Ulam AI / UnsolvedMath Contributors, CC BY 4.0. Source rights NEEDS\_REVIEW: publicly accessible author/publisher copy, redistribution terms requiring review.
\end{itemize}
\end{document}
